% ============================================================
%  Chapter 18 — Annihilating Polynomials and the Cayley–Hamilton Theorem
% ============================================================
\section{Annihilating Polynomials}
\label{sec:annihilating}

Given a linear operator $f$ on a finite-dimensional vector space, we can
substitute $f$ into any polynomial to obtain another linear operator.  The
polynomials that send $f$ to the zero operator form a rich algebraic
structure, culminating in the \emph{minimal polynomial} and the celebrated
Cayley--Hamilton theorem.

% ---- Polynomial of a linear operator ------------------------
\begin{defbox}[Definition --- Polynomial of a Linear Operator]
\label{def:poly-of-op}
Let $V$ be a finite-dimensional vector space over $\F$ and let $f$ be a
linear operator on $V$.  For a polynomial
$p(x)=\sum_{i=0}^{m}a_i x^i\in\F[x]$, define the linear operator
\[
    p(f) = \sum_{i=0}^{m} a_i\,f^i,
    \qquad f^0=\Id_V,\quad f^i = f\circ f^{i-1}\ (i\ge 1).
\]
\end{defbox}

\begin{defbox}[Definition --- Annihilating Polynomial and the Set $L_f$]
\label{def:annihilating}
Let $V$ be a finite-dimensional vector space over $\F$ and $f$ a linear
operator on $V$.  Define
\[
    L_f = \{p(x)\in\F[x]\mid p(f)=\mathrm{O}\}.
\]
A non-zero polynomial $p(x)\in L_f$ is called an
\textbf{annihilating polynomial} of $f$.

\medskip\noindent
The set $L_f$ satisfies:
\begin{enumerate}[label=\textbf{(\roman*)}]
    \item $0\in L_f$.
    \item $p,q\in L_f\Rightarrow p+q\in L_f$.
    \item $p\in L_f,\;h\in\F[x]\Rightarrow hp\in L_f$.
    \item $L_f\neq\F[x]$ (since $p(x)=1$ gives $p(f)=\Id_V\neq\mathrm{O}$).
\end{enumerate}
\end{defbox}

\begin{exbox}[Example --- Annihilating Polynomials for the Identity]
\label{ex:annihilating-id}
Let $V$ be a finite-dimensional vector space and $f=\mathrm{id}$.
\begin{itemize}
    \item $p(x)=2+6x+7x^2-9x^3$:
          $p(\mathrm{id})=(2+6+7-9)\mathrm{id}=6\,\mathrm{id}\neq\mathrm{O}$,
          so $p\notin L_{\mathrm{id}}$.
    \item $q(x)=x-1$:
          $q(\mathrm{id})=\mathrm{id}-\mathrm{id}=\mathrm{O}$,
          so $q\in L_{\mathrm{id}}$.
    \item $r(x)=x$:
          $r(\mathrm{id})=\mathrm{id}\neq\mathrm{O}$,
          so $r\notin L_{\mathrm{id}}$.
\end{itemize}
\end{exbox}

% ---- Existence of annihilating polynomial -------------------
\begin{thmbox}[Lemma --- Existence of an Annihilating Polynomial]
\label{lem:annihilating-exists}
Let $V$ be a finite-dimensional vector space over $\F$ and let $f$ be a
linear operator on $V$.  There exists a non-zero polynomial
$p(x)\in\F[x]$ such that $p(f)=\mathrm{O}$.
\end{thmbox}

\begin{tipbox}[Remark --- Bound on the Minimal Degree]
Since $L_f\neq\{0\}$, the number
\[
    l = \min\{\deg p(x) \mid p(x)\in L_f\setminus\{0\}\}
\]
is well-defined and satisfies $l\le n^2$ (where $n=\dime V$), because
$\dime_\F(\End(V))=n^2$, so the operators
$\Id_V, f, f^2,\ldots,f^{n^2}$ cannot all be linearly independent in $\End(V)$.
\end{tipbox}

\begin{propbox}[Proposition --- Existence and Uniqueness of the Monic Annihilating Polynomial of Minimal Degree]
\label{prop:minimal-poly-unique}
Let $V$ be a finite-dimensional vector space over $\F$ and let $f$ be a
linear operator on $V$.  There exists a unique \emph{monic} polynomial
$m(x)\in L_f$ of degree $l$ (the minimal degree defined above).
\end{propbox}

% ---- Minimal polynomial -------------------------------------
\begin{defbox}[Definition --- Minimal Polynomial]
\label{def:minimal-poly}
Let $V$ be a finite-dimensional vector space over $\F$ and $f$ a linear
operator on $V$.  The \textbf{minimal polynomial} of $f$, denoted
$m_f(x)$, is the unique monic polynomial $m(x)\in\F[x]$ satisfying:
\begin{enumerate}[label=\textbf{(\roman*)}]
    \item $m(x)$ is monic.
    \item $m(f)=\mathrm{O}$.
    \item $\deg m\le\deg p$ for all $p\in\F[x]$ with $p(f)=\mathrm{O}$.
\end{enumerate}
\end{defbox}

\begin{tipbox}[Remark --- Properties of the Minimal Polynomial]
\label{rem:min-poly-props}
\begin{enumerate}[label=\textbf{(\arabic*)}]
    \item The minimal polynomial $m_A(x)$ of a matrix $A\in\F^{n\times n}$
          is defined analogously.
    \item Similar matrices have the same minimal polynomial.
    \item If $A=[f]_{\mathcal{B}}$ for some basis $\mathcal{B}$, then
          $m_f(x)=m_A(x)$.
    \item $p(f)=\mathrm{O}$ if and only if $m_f(x)\mid p(x)$.
\end{enumerate}
\end{tipbox}

% ---- Examples of minimal polynomials -----------------------
\begin{exbox}[Example --- Minimal Polynomial of the Identity]
\label{ex:min-poly-id}
For $f=\mathrm{id}$, suppose $m_f(x)=x+a_0$ (testing degree 1).  Then
\[
    \mathrm{O}=m_f(f)=\mathrm{id}+a_0\,\mathrm{id}=(1+a_0)\,\mathrm{id}
    \;\Longrightarrow\; a_0=-1.
\]
Hence $m_f(x)=x-1$.
\end{exbox}

\begin{exbox}[Example --- Minimal Polynomial of a Rotation Matrix]
\label{ex:min-poly-rotation}
Let $A=\begin{bmatrix}0&-1\\1&0\end{bmatrix}$ (rotation by $90^\circ$).

\medskip\noindent
\textbf{Degree 1 impossible:} If $m_A(x)=x+a_0$, then
$A+a_0 I=\begin{bmatrix}a_0&-1\\1&a_0\end{bmatrix}\neq\mathrm{O}$.
So $\deg m_A\ge 2$.

\medskip\noindent
\textbf{Degree 2:}
$A^2=\begin{bmatrix}-1&0\\0&-1\end{bmatrix}=-I$,
so $A^2+I=\mathrm{O}$.
Hence $m_A(x)=x^2+1$.
\end{exbox}

% ---- Eigenvalue lemma ---------------------------------------
\begin{thmbox}[Lemma --- Eigenvalue and Polynomial of an Operator]
\label{lem:eigenvalue-poly}
Let $V$ be a finite-dimensional vector space over $\F$, let $f$ be a
linear operator on $V$, and let $\lambda\in\F$ and $v\in V$ satisfy
$f(v)=\lambda v$.  Then for all $P\in\F[x]$:
\[
    P(f)(v) = P(\lambda)\,v.
\]
\end{thmbox}

% ---- Same roots theorem -------------------------------------
\begin{thmbox}[Theorem --- Minimal and Characteristic Polynomials Have the Same Roots]
\label{thm:min-char-same-roots}
Let $f:V\to V$ be a linear operator on an $n$-dimensional vector space
$V$ over $\F$.  The minimal polynomial $m_f(x)$ and the characteristic
polynomial $\chi_f(x)$ have exactly the same roots (in the algebraic
closure of $\F$).
\end{thmbox}

% ---- Diagonalizable l.t. corollary -------------------------
\begin{corbox}[Corollary --- Minimal Polynomial of a Diagonalizable Operator]
\label{cor:min-poly-diag}
A linear operator $f$ is diagonalizable with distinct eigenvalues
$\lambda_1,\lambda_2,\ldots,\lambda_k$ if and only if
\[
    m_f(x) = (x-\lambda_1)(x-\lambda_2)\cdots(x-\lambda_k).
\]
\end{corbox}

\begin{corbox}[Corollary --- Inverse as a Polynomial of an Operator]
\label{cor:inverse-poly}
Let $V$ be a finite-dimensional vector space over $\F$ and $f:V\to V$ a
linear operator.  If $f$ is invertible, then there exists a polynomial
$p(x)\in\F[x]$ such that
\[
    f^{-1} = p(f).
\]
\end{corbox}

% ============================================================
%  Cayley–Hamilton Theorem
% ============================================================
\subsection{The Cayley--Hamilton Theorem}

\begin{thmbox}[Theorem --- Cayley--Hamilton Theorem]
\label{thm:cayley-hamilton}
Let $V$ be a finite-dimensional vector space over $\F$ and let $f$ be a
linear operator on $V$.  Then
\[
    m_f(x) \mid \chi_f(x),
    \qquad\text{equivalently,}\qquad
    \chi_f(f) = \mathrm{O}.
\]
\end{thmbox}

\begin{thmbox}[Theorem --- Diagonalizability and the Minimal Polynomial]
\label{thm:diag-min-poly}
Let $V$ be a finite-dimensional vector space over $\F$ and $f$ a linear
operator on $V$.  Then $f$ is diagonalizable if and only if
\[
    m_f(x) = \prod_{i=1}^k (x-\lambda_i),
\]
where $\lambda_1,\lambda_2,\ldots,\lambda_k$ are the \emph{distinct}
eigenvalues of $f$.
\end{thmbox}

\begin{exbox}[Example --- Non-Diagonalizable Matrix via the Minimal Polynomial]
\label{ex:non-diag-min-poly}
Let
\[
    A = \begin{bmatrix}
        1&0&0&0&0\\
        1&1&0&0&0\\
        0&0&-1&0&0\\
        0&0&0&-1&0\\
        0&0&0&0&3
    \end{bmatrix}.
\]
Then $\chi_A(x)=(x-1)^2(x+1)^2(x-3)$ and $m_A(x)=(x-1)^2(x+1)(x-3)$.

\smallskip
Since $m_A(x)$ contains the repeated factor $(x-1)^2$, by
Theorem~\ref{thm:diag-min-poly} the matrix $A$ is \textbf{not
diagonalizable}.  One can verify directly:
$(A-I)(A+I)(A-3I)\neq\mathrm{O}$, while
$(A-I)^2(A+I)(A-3I)=\mathrm{O}$.
\end{exbox}

% ---- Proof of Cayley-Hamilton --------------------------------
\subsection{Proof of the Cayley--Hamilton Theorem}

The full proof of the Cayley--Hamilton theorem proceeds via two lemmas.

\begin{thmbox}[Lemma 1 --- Cyclic Basis and Cayley--Hamilton]
\label{lem:cyclic-basis-ch}
Let $V$ be an $n$-dimensional vector space over $\F$ and $f$ a linear
operator on $V$.  If there exists a vector $v\in V$ such that
$\{v, f(v), f^2(v), \ldots, f^{n-1}(v)\}$ is a basis of $V$
(a \emph{cyclic basis} with generator $v$), then $\chi_f(f)=\mathrm{O}$.
\end{thmbox}

\begin{tipbox}[Remark --- Polynomial of a Block Upper-Triangular Matrix]
\label{rem:block-poly}
Let
$A=\begin{bmatrix}A_1&B\\0&A_2\end{bmatrix}\in\F^{n\times n}$
with $A_1\in\F^{n_1\times n_1}$, $A_2\in\F^{n_2\times n_2}$,
$n=n_1+n_2$.  For any polynomial $p(x)\in\F[x]$:
\[
    p(A) = \begin{bmatrix}p(A_1)&C\\0&p(A_2)\end{bmatrix}
\]
for some matrix $C\in\F^{n_1\times n_2}$.

\medskip\noindent
\emph{Verification.}  Writing $p(x)=a_0+a_1 x+\cdots+a_m x^m$:
\begin{align*}
    p(A)
    &= a_0 I_n + a_1 A + a_2 A^2 + \cdots + a_m A^m \\
    &= \begin{bmatrix}a_0 I_{n_1}&0\\0&a_0 I_{n_2}\end{bmatrix}
     + a_1\begin{bmatrix}A_1&B\\0&A_2\end{bmatrix}
     + a_2\begin{bmatrix}A_1^2&A_1 B+BA_2\\0&A_2^2\end{bmatrix}
     + \cdots
     + a_m\begin{bmatrix}A_1^m&*\\0&A_2^m\end{bmatrix} \\
    &= \begin{bmatrix}p(A_1)&C\\0&p(A_2)\end{bmatrix}.
\end{align*}
\end{tipbox}

\begin{thmbox}[Lemma 2 --- Block Matrix and Cayley--Hamilton]
\label{lem:block-ch}
Let $A=\begin{bmatrix}A_1&B\\0&A_2\end{bmatrix}\in\F^{n\times n}$ with
$A_1\in\F^{n_1\times n_1}$, $A_2\in\F^{n_2\times n_2}$, $n=n_1+n_2$.
If $\chi_{A_1}(A_1)=0$ and $\chi_{A_2}(A_2)=0$, then $\chi_A(A)=0$.
\end{thmbox}

\begin{proofbox}[Proof --- Cayley--Hamilton Theorem (Outline)]
The proof proceeds by induction on $n=\dime V$.

\textbf{Base case} $n=1$: $f=\lambda\,\mathrm{id}$, so $\chi_f(x)=x-\lambda$
and $\chi_f(f)=f-\lambda\,\mathrm{id}=\mathrm{O}$. \checkmark

\textbf{Inductive step}: If there exists a cyclic vector $v$ (i.e.\
$\{v,f(v),\ldots,f^{n-1}(v)\}$ is a basis), apply Lemma~\ref{lem:cyclic-basis-ch}.
Otherwise, there exists a proper $f$-invariant subspace $W\subsetneq V$, so
$[f]$ can be written in block upper-triangular form
$\begin{bmatrix}A_1&B\\0&A_2\end{bmatrix}$ with smaller diagonal blocks; apply
the inductive hypothesis to each block and then Lemma~\ref{lem:block-ch}.
\hfill$\blacksquare$
\end{proofbox}
